% Generated draft: requires mathematical and source review.
\documentclass[11pt]{article}
\usepackage[margin=25mm]{geometry}
\usepackage{fontspec}
\setmainfont{Latin Modern Roman}
\usepackage{amsmath,amssymb,mathtools,unicode-math}
\setmathfont{Latin Modern Math}
\usepackage{enumitem,needspace,xurl,hyperref}
\setlength{\parindent}{0pt}
\setlength{\parskip}{5pt}
\setlength{\emergencystretch}{3em}
\newcommand{\N}{\mathbb{N}}
\newcommand{\Z}{\mathbb{Z}}
\newcommand{\Q}{\mathbb{Q}}
\newcommand{\R}{\mathbb{R}}
\newcommand{\C}{\mathbb{C}}
\newcommand{\F}{\mathbb{F}}
\newcommand{\A}{\mathbb{A}}
\newcommand{\PP}{\mathbb{P}}
\newcommand{\E}{\mathbb{E}}
\newcommand{\eps}{\varepsilon}
\newcommand{\dd}{\,\mathrm{d}}
\begin{document}
\section*{UnsolvedMath ID: 30000140}
\textbf{Source problem code:} catalogue entry
\subsection{Definitions and mathematical statement}
For a positive integer \(n\), let \(\mathcal P(n)\) be the following assertion. There exist a finite constant \(D_n>0\) and a sequence of smooth, closed, simply connected \(n\)-manifolds with Riemannian metrics \((M_j,g_j)_{j\geq1}\) such that
\[
 M_i\simeq M_j\quad\text{for all }i,j,\qquad
 M_i\not\cong_{\mathrm{homeo}}M_j\quad\text{for }i\ne j,
\]
and, for every \(j\), every point \(x\in M_j\), every two-dimensional tangent plane \(\sigma\subset T_xM_j\), and every nonzero \(v\in T_xM_j\),
\[
 0\leq\operatorname{sec}_{g_j}(\sigma)\leq1,
 \qquad \operatorname{Ric}_{g_j}(v,v)>0,
 \qquad \operatorname{diam}(M_j,g_j)\leq D_n.
\]
Here closed means compact without boundary, \(\simeq\) denotes homotopy equivalence, and \(\cong_{\mathrm{homeo}}\) denotes homeomorphism of the underlying topological manifolds. Simple connectivity is understood to include connectedness. All the conditions concern the same chosen metric \(g_j\) on each manifold. The diameter bound is uniform in \(j\), after the sectional-curvature normalization to \([0,1]\). Positive Ricci curvature is pointwise positive definiteness; the statement does not impose a positive Ricci lower bound uniform in the sequence or a volume lower bound.

The cleaned question asks to determine
\[
 n_{\min}=\min\{n\in\Z_{>0}:\mathcal P(n)\}.
\]
Determining this minimum requires existence in the proposed dimension and exclusion of every smaller positive dimension. This notation presupposes that the set is nonempty; the supplied record does not reproduce an existence theorem or a candidate value.

The original statement asks for the smallest dimension with nonnegative sectional curvature, positive Ricci curvature and ``the conclusion of Theorem A.'' The verification field reports that the official report was checked and that the placeholder was replaced by the full curvature and diameter conditions of the preceding theorem. The formulation above follows that supplied cleaned version, including the infinite homotopy-equivalent, pairwise nonhomeomorphic family. The actual Theorem A is not reproduced, so equivalence with its full conclusion cannot be checked from this record alone.

The literature assessment checked on 2026-08-21 labels the problem open but describes the extracted sentence as a connective fragment, with a surrounding finiteness question and partial counterexamples. That description does not assess the fully expanded cleaned statement explicitly and is retained as a source-version qualification, rather than used to infer a dimension or a current result.
\subsection{Short English statement}
Find the smallest dimension supporting an infinite family of mutually homotopy-equivalent but pairwise nonhomeomorphic closed simply connected manifolds, each carrying a metric with sectional curvature in \([0,1]\), positive Ricci curvature and one common diameter bound. The cleaned statement expands a reference to Theorem A, whose full text is absent; the older literature note still describes an incomplete extraction.
\subsection{Sources}
\begin{itemize}
\item Local UnsolvedMath registry, permanent ID 30000140, source code catalogue entry, ``Least Dimension for Curvature-Finiteness Counterexamples,'' including original and cleaned statements and the supplied verification of the expansion of Theorem A.
\item \emph{Geometrie (2005)}, Oberwolfach Reports, pp.\ 2491--2538, DOI 10.4171/owr/2004/47: \url{https://doi.org/10.4171/owr/2004/47}. This is also the supplied literature source labelled ``original report.'' The citation says 2005, while the proposed-year metadata and DOI use 2004.
\item Literature triage checked 2026-08-21; local record updated 2026-09-16. The dated triage describes an incomplete extraction, whereas the verification field describes a later expansion. No cited source has been visited or independently checked here.
\end{itemize}
\end{document}
